\documentclass[10pt,a4paper]{article}
\usepackage[T1]{fontenc}
\usepackage[utf8]{inputenc}
\usepackage{lmodern,microtype}
\usepackage[margin=18mm]{geometry}
\usepackage{amsmath,booktabs,tabularx,array}
\usepackage{xcolor,hyperref,enumitem}
\definecolor{navy}{HTML}{17365D}
\hypersetup{colorlinks=true,linkcolor=navy,urlcolor=navy}
\setlength{\parindent}{0pt}
\setlength{\parskip}{0.55\baselineskip}
\setlist{nosep,leftmargin=*}
\newcolumntype{Y}{>{\raggedright\arraybackslash}X}
\newcommand{\code}[1]{\texttt{\detokenize{#1}}}
\title{\textbf{Time-series Foundation Model Evaluation}\\Experiment Guideline}
\author{Thesis project}
\date{27 August 2026}
\begin{document}
\maketitle

\section{Scope}
This is an inference-only evaluation of frozen Chronos-2, Chronos-Bolt, and
Chronos-T5, official TS-ICL, and four legacy TimeTensor
baselines: persistence, expected value, repeat, and weekly aligned lookback.
It reproduces and extends the accepted ICLR TSALM covariate study. It contains
no optimization, forecasting-model training, PatchTST training, or
data-generating process.

\section{Forecasting and evaluation contract}
Targets are wide univariate panels with shape
$(\text{users},1,\text{dates})$. For query date $t$,
\[
X=(t-L,t],\qquad Y=(t,t+H].
\]
Query dates are deterministic, date-major, spaced by
\code{evaluation.stride} (512 for the paper protocol), and evaluated for every
user. Evaluation fractions own target dates. Pairwise constant-window removal
removes only user/query pairs whose lookback standard deviation is at most the
configured threshold.

Target-only reversible instance normalization standardizes the target history
before inference and reverses it on predictions. Covariates remain on their
raw source scale. Evaluation nMSE always divides errors by the lookback scale,
independently of preprocessing.

\section{Models and covariates}
\begin{center}
\begin{tabularx}{\linewidth}{@{}lY@{}}
\toprule
Axis & Cases \\
\midrule
Models & Persistence, expected value, repeat, weekly aligned lookback, frozen
Chronos-2, Chronos-Bolt, Chronos-T5, and official TS-ICL. \\
Covariates & \code{none}; \code{identity} (target history supplied as known
covariate); \code{known} (prepared target-aligned covariate panels). \\
Controls & Instance normalization on/off; constant-window removal on/off. \\
\bottomrule
\end{tabularx}
\end{center}
External covariate panels must share the target timeline and either match user
names/order or contain one global column.

The lookback baseline infers the number of observations in seven days from
dated inputs, or accepts an explicit positive period override. For horizon
$H$ and period $P$, it returns the latest complete $H$-step history block whose
start is shifted back by $\lceil H/P\rceil P$ observations. This preserves
weekly alignment and prevents overlap with the forecast horizon. Inputs without
a cadence that divides one week require an explicit period. The former
absolute-index variants are not part of the current model set.

All foundation-model adapters use their official packages. TS-ICL retains
native preprocessing and test-time behavior and selects the median point
forecast. The foundation workflow supplies raw targets and disables the outer
instance-normalization wrapper. The sole shared aliases are
\code{chronos2}, \code{chronos_bolt}, \code{chronos_t5}, and
\code{ts_icl}. Covariate configuration does not alter model identity, and
an adapter without covariate support raises on a non-empty covariate input.
Expected checkpoints are
\begin{itemize}
  \item Chronos-2: \path{weights/chronos2};
  \item Chronos-Bolt: \path{weights/chronos-bolt-base};
  \item Chronos-T5: \path{weights/chronos-t5-base};
  \item TS-ICL: \path{weights/tsicl/tsicl-v1.ckpt};
\end{itemize}

\section{Metrics and seed}
Every run reports MSE, MAE, nMSE, nMAE, and MASE with configurable seasonality.
For each metric it records the sample mean and population standard deviation,
equal-user mean and population standard deviation across user means, the mean
of the worst $\lceil0.1U\rceil$ users, per-user within-user dispersion, and
per-horizon mean and dispersion. It also records total model inference time and
amortized time per user and per series/window; model and data loading are
excluded from inference timing. One run seed fixes every stochastic choice
exposed by a model or sampler; deterministic evaluation ordering remains fixed.

\section{Experiment families and grid}
The DGX fronts under \path{slurm/dgx/main/} are:
\begin{enumerate}
\item \code{01_univariate.slurm}: the four baselines and Chronos-2;
\item \code{02_controls.slurm}: normalization and constant-window controls;
\item \code{03_covariates.slurm}: univariate and identity-covariate inference,
      comparing Chronos-2 with TS-ICL, with prepared known covariates available
      by override;
\item \code{04_foundation_models.slurm}: frozen Chronos-2, Chronos-Bolt,
      Chronos-T5, and TS-ICL with model-native preprocessing.
\end{enumerate}
\code{test} is Electricity at $504{:}168$, seed 1. Full profiles use three
frequency-aware forecast ranges:
\begin{center}
\small
\begin{tabular}{lccc}
\toprule
Cadence & Short & Mid & Long \\
\midrule
Hourly & $168{:}24$ & $336{:}48$ & $504{:}168$ \\
Daily & $7{:}1$ & $14{:}2$ & $30{:}7$ \\
15-minute & $96{:}4$ & $192{:}8$ & $672{:}96$ \\
\bottomrule
\end{tabular}
\end{center}
The first three families apply all three ranges to Electricity, Traffic,
Solar, Exchange Rate, and every feasible dataset discovered in
\path{datasets/time/catalog.json}. TIME tasks shorter than $L+H$ are omitted.
\code{ultra} additionally selects Weather and ETTh1. Supersets select from the
same identity trees and reuse only signature-matched completed runs.

The foundation-model family uses Electricity at $504{:}168$ for \code{test};
its \code{full} profile applies the cadence-specific mid and long ranges to
Electricity, Traffic, Solar, Weather, Exchange Rate, and feasible TIME
datasets. Its report fixes Chronos-2 as the default reference. \code{ultra}
uses all three ranges and adds ETTh1.

\section{Execution}
\begin{verbatim}
EXPERIMENT_MODE=test sbatch slurm/dgx/main/01_univariate.slurm
EXPERIMENT_MODE=full sbatch slurm/dgx/main/01_univariate.slurm
EXPERIMENT_MODE=full sbatch slurm/dgx/main/02_controls.slurm
EXPERIMENT_MODE=full sbatch slurm/dgx/main/03_covariates.slurm
EXPERIMENT_MODE=test sbatch slurm/dgx/main/04_foundation_models.slurm
EXPERIMENT_MODE=full sbatch slurm/dgx/main/04_foundation_models.slurm
\end{verbatim}
Each command has a matching \code{_selena.slurm} overflow front under
\path{slurm/selena/main/}. It executes
the same scientific implementation on Selena but writes Slurm streams under
\path{logs_selena/} and all manifests, results, and reports under
\path{outputs_selena/}. DGX fronts retain \path{logs/} and \path{outputs/}.
Selena uses partition \code{an} with QoS \code{an_preemptable}.
The default workflow is \code{STAGES=evaluate,report}; stage selection is a
recovery override. Dataset, setting, model, seed, normalization, constant
handling, and covariates can be narrowed with the
documented environment overrides. Foundation-model inference is remote only.
Explicit dataset and setting overrides replace automatic selection. Without a
dataset override, full and ultra profiles require the TIME catalog; an explicit
\code{DATA_ROOT} names the datasets root containing \path{time/catalog.json}.
\code{LOOKBACK_PERIOD_STEPS} overrides weekly cadence inference for the
lookback baseline, and \code{MASE_SEASONALITY} sets the MASE denominator lag.

Slurm workflows never submit a publisher or execute Git commands. After any
terminal state, including failure, cancellation, or timeout, the user runs
\code{bash publish_job.sh <job-id>} from the project root. The script verifies
\code{main}, sources \code{$HOME/codes/proxy.sh} (or
\code{PROXY_SCRIPT_PATH}), and fast-forward pulls \code{origin/main} before
artifact selection or staging. Lightweight publication selects the exact
local or Selena stdout/stderr pair when a job ID is supplied, all log trees
otherwise, and only aggregate \path{outputs*/reports/} artifacts. Detailed
publication adds all non-binary outputs and diagnostics. Both tiers exclude
\code{*.pt}, \code{*.npy}, and \code{*.cbm}. The script then pushes
\code{origin/main} without opening a pull request. A
non-fast-forward pull stops without a merge commit, and unrelated staged paths
remain outside the publisher commit.

\section{Data provenance and artifacts}
The loader discovers \code{config.json} beside the selected CSV.
Portable settings are top-level, project overrides live under
\code{tsfm_evaluation}, explicit run values override both, and
\code{drop_users} uses replacement precedence: null inherits, an empty list
keeps every CSV user, and a nonempty override replaces the prior list. The
selected path and applied keys are logged. TSFM does not consume TimeTensors
tensor caches.

\subsection*{Selective TIME preparation}
\code{python -m scripts.prepare_time_csv} resolves an immutable revision of
\path{Real-TSF/TIME-ProcessedCSV} and downloads only selected CSV paths. The
default considers 15-minute, hourly, and daily sources, retaining source tasks
with at most 500 usable series and at most 10,000 dates per series. Files are
merged only when their complete timestamp indices match; non-aligned groups
become separate datasets, while irregular, short, or non-finite series are
recorded as skips. Each accepted panel is written below
\path{datasets/time/<dataset>/} with a directly loadable wide CSV and
\path{config.json}. \path{datasets/time/catalog.json} records source revision,
counts, skips, and eligible query/window totals. By default, each source uses
only the hourly, daily, or 15-minute short/mid/long ranges in the table above;
an explicit settings option applies a common custom grid. The current revision
inventory retains five hourly and seven daily
source configurations under the default limits; no 15-minute source qualifies.

The four independent family roots live below \path{outputs/} on DGX and below
\path{outputs_selena/} on Selena. Their relative names are
\code{univariate}, \code{controls}, \code{covariates}, and
\code{foundation_models}. Their ordered identity is
\begin{verbatim}
outputs/family/dataset/L_H/backbone/covariate_mode/normalization/
  constant_policy/run_n/
\end{verbatim}
Every model config owns one directory. Batch size, data split and stride,
metric/window export, model inference options, and explicit covariate bindings
are pipeline configs; device and Slurm placement are runtime-only. One seed
fixes every stochastic choice. Mode filters path subsets and is not part of the
path or computation signature.
TSFM launches one seed per run, records it in the manifest, and stores its
artifacts directly under \path{run_n}; a seed leaf is reserved for a future
multi-seed launcher. The manifest structure, statuses, collision and selection
policies, purposes, and report manifests are thesis-wide conventions. The
workflow roots, ordered identity, and required summary artifacts are
TSFM-specific.

Each completed run contains \path{summary.json},
\path{per_user_metrics.csv}, \path{horizon_metrics.csv}, optional
\path{window_metrics.csv}, \path{resolved_config.json}, and a
\code{schema_version: 1} manifest with
declared configs, optional plain provenance, launch attempts, required artifacts,
per-seed state, and one of four overall statuses; seed state may additionally be
\code{ready}. Run identity contains only the schema,
identity/model configs, pipeline/experiment configs, and seeds. It never
fingerprints or hashes code, Slurm files, data, weights, logs, outputs, or directories;
code/data changes are manual rerun decisions. The schema changes only for a
deliberate global contract break, and \code{new} creates another repeat with
unchanged parameters. The default \code{overwrite_exact} policy skips exact
completed work, resumes exact interruption, and creates a new \code{run_n} for
a changed pipeline. The evaluator records ready artifacts, and completion is
written immediately after that evaluation's \code{srun} returns successfully.
Later evaluation or report failures preserve completed runs and interrupt only
unfinished work. Completed
manifests are authoritative, without later file hashing or synchronization
checks. Reports accept distinct/latest/average config handling,
selected/latest/distinct/average repeat selection, explicit pipeline filters,
and purpose filters. Explicit pipeline filters select configurations and must
match even with one run. Nested pipeline and experiment fields, including the
scientific identity embedded for an upstream dependency, are exposed through
dotted filter keys and every differing nested field appears in a distinct
label. If latest-run timestamps tie, the larger \code{run_n} index wins, making
selection deterministic. Reports write requested filters and obtained inputs to
\path{outputs/reports/family/mode} on DGX or the same relative path below
\path{outputs_selena/} on Selena, while \code{SELECTED_RUNS.txt} stores only
automatic or pinned exact repeats. Per-configuration plots and aligned inputs
created for averaging live under \path{outputs/diagnostics/family/mode}; the
report's plot index stores paths relative to that diagnostics directory.

An average policy aligns and averages the selected summary, per-window,
per-user, and per-horizon metric artifacts before computing comparisons,
Chronos-2 win rates, marginal tables, or plots. Thus every downstream analysis
uses the same averaged population rather than an arbitrary selected repeat.

Reports compare complete configurations without mixing incompatible pipeline
settings. They write the selected raw results, model-versus-reference
configuration table, equal-configuration marginal tables by dataset, by
cadence-independent short/mid/long range, and by literal setting, Chronos-2
strict paired-window win rates against the best aggregate
baseline, and a report manifest. The selected table metric may be any error
metric or total/per-user/per-window inference time. When window artifacts are
enabled, the diagnostics tree also receives per-user mean-versus-dispersion
scatters, window error histograms, and horizon-error curves as PNG and PDF,
while the compact report retains their plot index.

\vspace{0.5\baselineskip}
\section{Validation}
The shared scientific runtime may run dataset, metric, model-adapter,
repeat/report, manifest, and Slurm workflow checks. Local checks may use fake
adapters or the repeat baseline, but must not load Chronos-2, Chronos-Bolt,
Chronos-T5, or TS-ICL checkpoints or run full sweeps.
\end{document}
